\section{TedX Review}
\label{research:ReviewTedX}
Following on from my HPQ Project I had the incredible opportunity to present a TedX talk at my school. In preparing for this talk I had many more thoughts and many more notes written which i also want to reflect on.

\begin{center}
    \textbf{TedX Will Artificial Intelligence Render Our Workforce Obsolete?}\\
    By Henry Field \autocite{tedxtalksAIComingYour2024}
\end{center}

\sloppy
The Machines will throw thousands of your petitioners out of employment, whereby they are brought into great distress, are not able to procure a maintenance for their families, and are deprived of the opportunity to bring up their children to the workforce.

could this relate to AI? Given the current furore about the impact that it could have on jobs, life and even the continued existence of this planet!
Yet this quote originated not from this year, nor this decade or even this century but from 1786 and yet it perfectly encapsulates the situation we now see the workforce facing today . The 18th century brought about mass change as the industrial revolution sparked modernisation efforts across western Europe. Despite mass concern for workers welfare and employment, our population still stands strong and has since faced not 1, not 2, but 3 events that have been described as revolutions, so are we finally facing into an epoch moment or is this just another boy crying wolf . Today I aim to break down the recent scientific breakthroughs and media storm that has provoked this most recent time of uncertainty and hopefully reduce any fears which recent links with popular fiction may have given you.
perhaps the most used marketing buzz phrase of the year, The term AI can now been seen plastered against any number of digital services despite there being no clear links between their function and any common definition of Artificial intelligence at all! So what exactly Is AI.

As defined by the Organisation for Economic Co-operation and Development (OECD), AI is “a machine-based system that is capable of influencing Its environment by producing an output for a given set of objectives.” (OECD, 2019), hardly trips of the tongue does it

Britannica states AI as being “the ability of a digital computer or computer-controlled robot to perform tasks commonly associated with intelligent beings.” (Britannica, 2023) Despite this definition not being clear and changing not even with the seasons, but with the days of the week. The only agreeable definition is that information goes in to a magic box and a transformational output or work emerges. Now hold that thought.

Fortunately, a worker on the other hand is easy to define in both UK and international law. Defined as a person with a, “contract or other arrangement to do work or services in exchange for a reward” (UK Government, 2023). Once again this theme of output or work is found, though a new variable is added into the equation, namely - reward. Simply put, both AI and workers complete work but ONLY workers require reward - but lets put that aside for a minute.
It may seem that this technology has only just been unleashed - with a series of breakthroughs in quick succession .whereas, In actuality, development has been underway for over 70 years.

Alan Turing spoke of Artificial Intelligence in the 1950s proposing it In his paper, “Computing Machinery and Intelligence”. This introduced the now well known -Turing Test. A conceptual examination aimed to demonstrate a machine so adapt at reproducing human behaviour that it could engage in conversation with a human participant without their realisation of its own computational origins. No functional device was built at that point but it laid the foundations for a pathway that led us to future developments.

Lets introduce another key term that is used alongside artificial Intelligence, Machine learning. Independently designed and built in 1950, by researcher Claude Shannon was Theseus, a primitive mouse robot that could navigate a maze via trial and error while outputting its “memories” to a circuit, enabling it instantly solve the challenge In the future. Once again demonstrating input, process and output.

Machines such as Theseus utilised this machine learning technology to help a machine refine its own processes and produce better outputs when given the opportunity to learn from new data.

This definition does cause some people to form a link between this concept and the biological functioning of human learning, and at very high level it does both in process and scientific foundation.

From birth, humans learn to form their behaviours and knowledge of this world by learning and processing experiences and shared interactions. As children gradually increase their experience they grow into fully independent beings which we call adults.
To many this would make Machine learning appear equal in scope to a humans ability to learn but they would be wrong. the topic of this talk - Ai - bridges a crucial diffrence in this comparison.

Computers have long been seen as superior to humans when it comes to performing numerical or logical calculations. This is after all, what has led to their widespread adoption. Machine learning then furthers this by introducing new processing methods and tools for large data-sets. Computers however, do also have their weaknesses which can primarily be found in the difficulty to decipher language, emotion and less rational thought.

This is where AI intervenes! and becomes an additional layer on top which aims to merge many machine learning algorithms together into a single unified system that can effectively mimic a human's thought process. This biomimicry enables modern technologies such as ChatGPT to be so user-friendly as the technology is aiming to function as a human rather than just producing a human output. It can therefore seem to many that Artificial intelligence is a confusing and vague term that can only be described in relation to organic thought. Lets once again, put this thought to a side.

But what has made the creation of AI possible since models like the Theseus Machine. Scientists such as Arnaud Miribel see AI as having four fundamental building blocks which decide their effectiveness: Quantity of data, Machine learning, Computing power, and Human expertise (Miribel, 2023) - and that once humanity can combine these aspects, we will have a system sufficiently advanced to be defined as true AI - so let's break these down:

Data is said to be more valuable than money by some which has caused massive organisations (such as Meta, the parent company of Facebook) to release free products that sustain themselves via the collection and sale of this precious resource. Around 120 zettabytes of data are estimated to be transferred yearly and these statistics are only increasing. Google estimates that only around 1.2million data points are required for Ai to function at a baseline and so I believe it can be safely said that we have developed the capacity to form the large data sets required to train AI models.

Machine learning as previously mentioned has continued to grow as models such as AlphaGo by Google have begun dominating gaming circuits. Despite the rules of Go being relatively straightforward, sequentially adding stones to a 19x19 grid to gain territory, the average game lasts for 150 moves and cumulates to $10^360$ board configurations - vastly more possible results than atoms in the universe. This means the traditional algorithms that chess "AI", for a relatively simpler game have relied on cannot be used as the optimal move processing time would last many years. These massive achievements showcase the level of technology now available to our population and highlights AI's ability to quickly grow and improve.

Our capacity to produce computational power has also massively increased, transistor sizes (a common factor measured for the performance of a processor) which started at 5 micrometres may have seemed small in the 1980s but are now hugely bloated when compared to the cutting edge 0.34 nanometres chiplets in production - a 99.9932\% reduction in size.

As the final factor, Human expertise is a rather difficult philosophical principle to tackle as we as a population are still not sure what this entails however this has been simplified by scientists to only include the ways we as organic beings solve challenges and interact with one another. This includes the processing of “natural language”, which entails how humans would normally communicate with one another in different situations such as verbal converse or written text rather than the strictly structured languages such as Python or C which computers have been built to understand. Models that have successfully incorporated these techniques have only recently started to appear in the public domain and have already revolutionised the way many see computers.

Now, over the course of me standing here we have now collected a series of 3 key points which I believe will define the future of our relationship with AI: Magic transformational outputs, lack of reward and comparisons to organic work, but how do these link with our predicted future.

In the past ten years, many startups have formed focused on improving AI technologies. Open AI is one such competitor, with their primary product, ChatGPT. This is one of the first examples of the technology that has become freely available to the general population and has shown how powerful AI can be, being capable of producing written reports, rewriting essays, and forming complex programs based solely on a written prompt. Despite this apparent success these machines are very much in their beta stages and have a habit of being confidently wrong - a large concern for the uptake of AI technologies. Google's project Bard - for example, has notoriously and repetitively struggled to answer the question of the first photo taken of a planet outside of our solar system - misidentifying it to be the James Webb Space Telescope rather than the correct, European Southern Observatory's VLT. The technology is not yet perfect and it is not possible to tell how long It will take for these machines to be reliable enough for constant use. Despite this, many companies are faced with fear of missing out and have already rushed to begin working on their implementation of AI within their products and workflows, a strategy that we will only be able to judge with hindsight.

We now have the benefits, risks and most importantly ability of this new technology defined, but how can we possible predict the magnitude of impact it will have in the short and long term future? Possibly by looking to the past. As I previously mentioned, although AI has not previously impacted our population, many other industrial events have throughout history. Seen as one of the earliest accounts of a mechanical boom, the 1760s brought about massive change in Britain as the Industrial Revolution shook the growing populace. Manual labour which had employed large swathes of the workforce was now being replaced by steam and ion, the steam engine revolutionised transport, heavily reduced the reliance on horse-drawn carriage and led to trains now carrying 100 passengers whilst requiring only 3 human operators, leading to massive cost savings. Many predicted that these changes would soon destroy the workforce, quickly enveloping every trade but these ideologies never came to fruition due to a singular inability of machinery - the ability to independently think.

The engine excelled at its singular job of providing drive however still required human operators to direct and fuel it while the loom required human ingenuity to set up and operate. These limitations led to entirely new industries of operations and maintenance being created - massively reducing the job-seeking population and enabling citizens who are challenged by physical disability to actively partake in the workforce. This shifting of common trades from that of labourers to trained operators allowed these new tools to become aids to amplify the skill and ability of any human.

Despite this positive precedent of greater equality and opportunity in the workforce, It has been predicted by Teamstage that 1.7 Million manufacturing jobs have been lost to automation since and that by 2030, 20 million manufacturing jobs could be lost. As companies have scaled, so have their mechanical operations and these machines have slowly consumed many of the previously accessible jobs they had originally created.

A recent report by Adecco Group found that only 4\% of workers believe their job will imminently be replaced by artificial intelligence. This statistic concerns many as the same study points to this figure being around 85\% and points to a lack of knowledge from the public at large. AI could become a massive aid to human workers if they understand how to use it as a tool to further upskill themselves but many are concerned that this necessity for education that may widen the pre-existing knowledge gap and cause further inequality. The UK Government's report on the impact of AI on jobs stated that those with only GCSE qualifications will see a 7\% further reduction in employment in the next 5 years and with recent strike action from numerous unions and calls for higher salaries, employers are looking for AI to fill vacant positions and help decrease running costs. 63\% of surveyed organisations stated that pressure to reduce costs will require the use of AI In their industries and a further 84\% said that AI would allow them to obtain or sustain a competitive advantage.

So In some form of conclusion, I believe that the future of AI in the workforce is not all doom and gloom, but that Governments need to cooperate to guide legislation in ensuring a safe and sustainable introduction of the technologies and crucially provide support and education to their population to help them to adapt to the world that will soon be upon them. With many new jobs being created, education systems will have to move away from withholding these advancements from students to incorporating them in learning as an aid similarly to how calculators have changed the field of mathematics. As education catches up to these modern innovations - employment requiring these skills will start to increase and along with this can come great success for you, our population and workforce. To quote Yann LeCun, “our intelligence is what makes us human, and AI is an extension of that quality” (LeCun, 2022)

\fussy
